% Diagram automatu stanow -- TikZ
\begin{tikzpicture}[
    >=Stealth,
    node distance=3.6cm,
    st/.style={circle,draw,thick,minimum size=1.15cm,inner sep=1pt,font=\sffamily\small},
    ev/.style={font=\small,fill=white,inner sep=1.5pt}
]
    \node[st] (S0)  at (0,0)    {$S_0$};
    \node[st] (S1)  at (5,1.9)  {$S_1$};
    \node[st] (S2)  at (5,-1.9) {$S_2$};
    \node[st] (S3a) at (10.5,1.9)  {$S_{3a}$};
    \node[st] (S3b) at (10.5,-1.9) {$S_{3b}$};

    % telefon wylaczony
    \path[->,thick]
        (S0) edge[loop above] node[ev] {K, Ś} (S0)
        (S0) edge[bend right=18] node[ev] {D} (S2);

    % przelaczanie ekranu
    \path[->,thick]
        (S1) edge[bend right=18] node[ev] {K} (S2)
        (S2) edge[bend right=18] node[ev] {K} (S1);

    % otwarcie i anulowanie menu
    \path[->,thick]
        (S1)  edge[bend left=18]  node[ev] {Ś} (S3a)
        (S3a) edge[bend left=18]  node[ev] {K} (S1)
        (S2)  edge[bend right=18] node[ev] {Ś} (S3b)
        (S3b) edge[bend right=18] node[ev] {K} (S2)
        (S3a) edge[loop right] node[ev] {Ś} (S3a)
        (S3b) edge[loop right] node[ev] {Ś} (S3b);

    % D z kazdego stanu z wlaczonym telefonem
    \node[draw,dashed,rounded corners,inner sep=0.45cm,
          fit=(S1)(S2)(S3a)(S3b)] (grp) {};
    \node[above=1pt of grp.north west,anchor=south west,font=\small\itshape]
        {stany z włączonym telefonem};
    \draw[->,thick] (grp.west |- S0) -- node[ev,pos=0.5] {D} (S0.east);
\end{tikzpicture}
